\chapter{FX market making}\label{ch:hf:fx-market-making}

Un fournisseur de liquidité non bancaire diffuse en continu des prix euro-dollar vers une centaine de banques, de courtiers et de plateformes à la fois, chacun avec son propre prix~; une transaction conclue sur
l'une d'entre elles est couverte, ou non, en quelques microsecondes sur une autre. Dans le modèle de ce chapitre, le last look du flux fait passer la capture sur le
palier le plus toxique de 0.13 à 0.20 point de base par demande, et la réduit sur les autres, pour un tiers de l'ensemble des demandes rejetées~; et internaliser le flux plutôt
que de le couvrir double le résultat de la journée, mais seulement si le skew du flux ramène les transactions de compensation dans les quelques centaines de transactions suivantes.

\section{Fourniture électronique de liquidité en principal}

Le marché des changes (le livre~2, chapitres~14 et~15) se négocie de gré à gré, dans un réseau de plateformes~: plateformes principales interdealer, réseaux de
communication électronique (ECN), plateformes single-dealer et multi-dealer, et les agrégateurs de liquidité du livre~10, chapitre~21. Les \omterm{def:hf:the-business-of-market-making:emm}{teneurs de marché
électroniques}, banques et non-bancaires, cotent en principal~: ils diffusent des prix, prennent l'autre côté de chaque transaction qu'ils acceptent, et gèrent l'inventaire
qui en résulte.

\begin{dated}{2026-09}{La taille du marché}\label{dat:hf:fx-market-making:bis}
Les résultats préliminaires de la BIS Triennial Central Bank Survey, publiés le 30 septembre 2025, situent les transactions de change à 9.6~billions de dollars par jour en avril 2025.
Les swaps de change restent l'instrument le plus échangé, à 4~billions de dollars par jour~; depuis avril 2022, le volume au comptant a progressé de 42~\% et celui des changes à terme sec (outright) de 60~\%, pour représenter 31~\%
et 19~\% du total. Le dollar américain figurait d'un côté dans 89~\% de l'ensemble des transactions~; les desks de vente du Royaume-Uni, des États-Unis, de Singapour et de Hong
Kong représentaient 75~\%.
\end{dated}

Chaboud, Chiquoine, Hjalmarsson et Vega ont étudié la montée du trading algorithmique sur le marché interdealer des changes, à l'aide de données qui identifiaient
l'activité générée par ordinateur. Le trading algorithmique a amélioré deux mesures de l'efficience des prix, la fréquence des opportunités d'arbitrage triangulaire et
l'autocorrélation des rendements à haute fréquence~; la baisse des opportunités d'arbitrage tient surtout aux ordinateurs qui prennent la liquidité, ce qui peut imposer des coûts
de sélection adverse plus élevés aux traders plus lents, et une autocorrélation moindre vient des ordinateurs qui en fournissent.

\section{Tarification à partir de plusieurs plateformes}

Le \omterm{def:hf:fair-value:fair}{prix juste} du fournisseur (chapitre~2) combine les plateformes qu'il observe, chacune avec son propre bruit et son propre délai. En pondérant le mid de chaque plateforme par
l'inverse de sa variance, le bruit plus la variance que le taux a accumulée depuis la dernière mise à jour de cette plateforme
(le \cref{lst:hf:fx-market-making:fair}), bat la plateforme principale seule~: dans le modèle du chapitre à trois plateformes (bruit de 0.3, 0.6 et 0.8 point de base,
se mettant à jour toutes les 1, 3 et 5 millisecondes en moyenne), le prix consolidé se situe à 0.26 point de base du taux vrai, contre 0.30 pour la plateforme principale.

\begin{definition}[Croisé synthétique]\label{def:hf:fx-market-making:cross}
Un \emph{croisé synthétique}\index{croisé synthétique} (synthetic cross) est un prix pour une paire de devises construit à partir de deux paires qui partagent une devise, comme l'euro-yen à partir
de l'euro-dollar et du dollar-yen, avec pour cours acheteur le prix des jambes auquel le fournisseur peut vendre, et pour cours vendeur celui auquel il peut acheter.
\end{definition}

Avec l'euro-dollar à 1.0850/1.0851 et le dollar-yen à 150.20/150.21, l'euro-yen synthétique vaut 162.967/162.993, soit une largeur de 2.6 pips, les spreads des deux jambes
se cumulant. Un fournisseur qui cote le cours croisé plus serré que son synthétique doit s'attendre à internaliser une partie du flux, ou à couvrir
sur le marché du cours croisé direct lorsque celui-ci est moins cher.

\section{Flux, paliers et last look}

Le fournisseur diffuse un prix à chaque client ou plateforme, plus large pour les paliers dont le flux est plus informé (le livre~9, chapitre~24,
sur la segmentation des clients (tiering)). Le last look (le livre~2, chapitre~15, où sont exposés les principes du FX Global Code à ce sujet) lui permet
de retenir une demande pendant quelques millisecondes et de la rejeter si le prix a bougé. Les flux du chapitre~: palier 1 (banques et plateformes, avec des clients d'arbitrage rapides) à 0.3 point de
base, dont 30~\% de ses demandes sont informées~; palier 2 (courtiers) à 0.5, dont 10~\% informées~; palier 3 (entreprises) à 1.0, aucune. Une demande informée suit un
mouvement de 0.6 point de base que le flux n'a pas encore affiché. Retenir pendant 5 millisecondes~; rejeter de façon asymétrique toute demande dont le prix a bougé de plus de 0.1 point
de base en faveur du client, ou de façon symétrique dans les deux sens (le \cref{lst:hf:fx-market-making:stream}).

\begin{omfigure}
\begin{tikzpicture}[font=\footnotesize,>=stealth,box/.style={draw,rounded corners,fill=omThm!8,minimum height=0.7cm,align=center}]
\draw[->,thick] (0,0) -- (11.2,0) node[right] {temps};
\foreach \x/\l in {0.6/{le flux affiche\\$S$},3.2/{la demande à $S$\\arrive},7.6/{contrôle après\\5 ms}} \draw (\x,0.12) -- (\x,-0.12) node[below,align=center] {\l};
\draw[decorate,decoration={brace,amplitude=5pt}] (3.2,0.35) -- node[above=4pt] {fenêtre de rétention} (7.6,0.35);
\draw[omDef,thick] (0.6,1.3) -- (2.4,1.3) -- (2.4,1.9) -- (7.6,1.9);
\node[font=\scriptsize,anchor=west,omDef] at (2.45,1.6) {le taux vrai bouge de 0.6 pb};
\node[font=\scriptsize,anchor=east] at (0.55,1.3) {taux};
\node[box] (acc) at (10.1,1.9) {accepter~: exécution\\au prix $S$};
\node[box,fill=omDef!15] (rej) at (10.1,0.95) {rejeter~: prix bougé\\$>0.1$ pb};
\draw[->] (7.7,1.3) -- (acc.west); \draw[->] (7.7,1.3) -- (rej.west);
\node[font=\scriptsize,align=left,anchor=north west] at (0.2,-1.05) {asymétrique~: on ne rejette que les mouvements en faveur du client (ceux de la demande informée)~;\\
symétrique~: on rejette les mouvements dans les deux sens, y compris ceux en faveur du fournisseur.};
\end{tikzpicture}

\omcaption{Last look tel que le flux du chapitre l'applique~: le fournisseur affiche un prix $S$~; la demande d'un client au prix $S$ arrive, les demandes informées après un
mouvement de 0.6 point de base du taux vrai que le flux n'a pas encore affiché~; le fournisseur la retient 5 millisecondes et la rejette si le prix a bougé
de plus de 0.1 point de base, seulement en faveur du client (asymétrique) ou dans les deux sens (symétrique). Source~:
\texttt{firm.fxlp} sur \texttt{firm.lastlook}.}\label{fig:hf:fx-market-making:lastlook}
\end{omfigure}

\begin{center}
\small
\begin{tabular}{@{}lrrrrrr@{}}
\toprule
& \multicolumn{2}{c}{sans last look} & \multicolumn{2}{c}{asymétrique} & \multicolumn{2}{c}{symétrique}\\
\cmidrule(lr){2-3}\cmidrule(lr){4-5}\cmidrule(l){6-7}
palier & capture & rejets & capture & rejets & capture & rejets\\
\midrule
palier 1 (30~\% informés) & 0.125 & 0 & 0.200 & 43.5~\% & 0.140 & 56.8~\%\\
palier 2 (10~\% informés) & 0.453 & 0 & 0.395 & 26.5~\% & 0.280 & 43.1~\%\\
palier 3 (aucun) & 1.008 & 0 & 0.851 & 18.3~\% & 0.637 & 37.0~\%\\
ensemble & 0.400 & 0 & 0.389 & 33.4~\% & 0.282 & 48.8~\%\\
\bottomrule
\end{tabular}
\end{center}

La capture est le mark-out du fournisseur par demande, exécutée ou non, en points de base (le mécanisme dans la
\cref{fig:hf:fx-market-making:lastlook}). Le last look ne rapporte que sur le palier toxique~: il rejette toutes les demandes informées
qui s'y présentent, et 19~\% des demandes non informées avec elles. Sur les autres paliers, il abandonne plus de bonnes transactions rejetées qu'il n'en économise, et la vérification
symétrique, qui rejette aussi les mouvements en faveur du fournisseur, abandonne encore davantage. Ce que le last look coûte aux clients (les rejets et
l'information qu'ils contiennent) relève du principe du FX Global Code à ce sujet~; l'arithmétique présentée ici est celle du fournisseur.

\section{Internaliser ou couvrir}

Chaque transaction acceptée alimente l'inventaire du fournisseur. Il peut couvrir immédiatement sur une plateforme principale, en payant le spread de cette plateforme, ou internaliser~: garder la
position et laisser les transactions clientes ultérieures la compenser, en réglant le skew de son flux pour que les clients soient plus enclins à prendre le côté qui la réduit. La
journée du chapitre (le \cref{lst:hf:fx-market-making:internalise})~: 20\,000 transactions clientes de 1~million de dollars, 0.4 point de base gagné sur chacune, une couverture à 0.25
point de base, le taux évoluant de 60 points de base sur la journée, et une pénalité sur la variance du P\&L d'inventaire. Le skew fait revenir l'inventaire
vers sa moyenne~; sa demi-vie, en transactions clientes, décide de la réponse.

\begin{omfigure}
\begin{tikzpicture}
\begin{axis}[width=11cm,height=5.4cm,axis y line*=right,axis x line=none,xmode=log,log basis x=10,xmin=50,xmax=2500,ymin=-250,ymax=800,
  ylabel={\small objectif (milliers de dollars par jour)},tick label style={font=\footnotesize},table/col sep=comma]
\addplot[omDef,thick,dashed,mark=square*,mark size=1.3pt] table[x=halflife,y=internalise]{figdata/hft/21-fx-market-making/hedge.csv};
\label{plot:hf:fxint}
\addplot[gray,thin] table[x=halflife,y=hedge]{figdata/hft/21-fx-market-making/hedge.csv};
\label{plot:hf:fxhedge}
\end{axis}
\begin{axis}[width=11cm,height=5.4cm,axis y line*=left,xmode=log,log basis x=10,xmin=50,xmax=2500,xtick={50,100,200,500,1000,2000},xticklabels={50,100,200,500,1000,2000},xlabel={\small demi-vie de l'inventaire (en
  transactions clientes)},ylabel={\small meilleure part couverte (\%)},ymin=0,ymax=105,tick label style={font=\footnotesize},grid=major,grid style={gray!25},
  table/col sep=comma,legend style={font=\scriptsize,at={(0.5,-0.3)},anchor=north},legend columns=3]
\addplot[omThm,thick,mark=*,mark size=1.3pt] table[x=halflife,y=best]{figdata/hft/21-fx-market-making/hedge.csv};
\addlegendentry{meilleure part de couverture (gauche)}
\addlegendimage{/pgfplots/refstyle=plot:hf:fxint}\addlegendentry{tout internaliser (droite)}
\addlegendimage{/pgfplots/refstyle=plot:hf:fxhedge}\addlegendentry{tout couvrir (droite)}
\end{axis}
\end{tikzpicture}

\omcaption{La part de chaque transaction client la mieux couverte immédiatement sur la plateforme principale (gauche), et l'objectif de la journée (fourchette moins les coûts de couverture moins
une charge de 0.002 par centaine de dollars au carré de variance du P\&L d'inventaire, droite) lorsque l'on internalise tout ou que l'on couvre tout, en fonction
de la demi-vie de l'inventaire sous le skew du flux (échelle logarithmique)~; 20\,000 transactions de 1~million de dollars, 0.4 point de base gagné, 0.25 payé pour couvrir. Données~:
\texttt{hf\_fx.hedging}.}\label{fig:hf:fx-market-making:hedge}
\end{omfigure}

Sans skew, l'inventaire est une marche aléatoire et la meilleure politique couvre 95~\% de chaque transaction. Avec une demi-vie d'environ 1\,150 transactions ou plus, elle couvre encore
95~\%~; à 870, 80~\%~; à 770, la moitié~; à 690 ou moins, presque rien (\cref{fig:hf:fx-market-making:hedge}). Dès que le flux se compense assez vite,
internaliser rapporte bien plus que couvrir~: 588\,000~\$ par jour pour une demi-vie de 350 transactions, contre 300\,000~\$ pour couvrir tout. La bascule
est brutale, et la raison mérite d'être connue. Sans skew, la variance de l'inventaire croît comme le carré de la part non couverte $1-h$, et l'objectif
$S-hC-\lambda(1-h)^2V$ atteint son maximum en $h^\ast=1-C/(2\lambda V)$~: 97~\% ici. Avec le skew, un inventaire plus important est ramené plus fermement, et la
variance ne croît qu'à peu près proportionnellement à $1-h$~: l'objectif est alors presque linéaire en $h$, et son maximum se situe à un coin.

\section{Croisés et marchés émergents}

Les croisés et les devises des marchés émergents sont là où l'internalisation et la synthèse comptent le plus~: leurs marchés directs sont peu profonds, leurs branches
passant par le dollar sont liquides, et un fournisseur qui reçoit du flux sur de nombreuses paires peut compenser l'euro-yen d'un client avec l'euro-dollar d'un autre et le dollar-yen d'un troisième.
Les contrats à terme non livrables (livre~2, chapitre~18) étendent la \omterm{def:hf:the-business-of-market-making:emm}{teneur de marché électronique} aux devises soumises à un contrôle des capitaux~; ils se règlent en dollars sur un
fixing, de sorte que le risque du fournisseur court jusqu'à la date de fixing, et sa couverture est un autre contrat à terme non livrable ou le marché onshore auquel il a accès.

\section{Fiches de stratégie}

\begin{strategyfile}{Teneur de marché de change sur la plateforme principale}\label{strat:hf:fx-market-making:primary}
\sfield{Qui vous paie, et pourquoi} Les banques et les fonds qui prennent la liquidité sur les plateformes interbancaires.
\sfield{Instruments et plateformes} Les paires majeures sur les plateformes interbancaires principales.
\sfield{Signal} Un \omterm{def:hf:fair-value:consolidated}{prix juste consolidé} à partir de toutes les plateformes, pondéré par le bruit et l'obsolescence.
\sfield{Dimensionnement et exécution} Coter au meilleur cours avec un skew d'inventaire~; cotations fermes, sans last look.
\sfield{Coûts} Frais de plateforme~; sélection adverse de la part des preneurs plus rapides.
\sfield{Comment elle disparaît} La vitesse~: les preneurs algorithmiques suppriment les opportunités d'arbitrage et imposent une sélection adverse aux fournisseurs plus lents.
\sfield{Horizon, capacité, infrastructure} De la microseconde à la seconde~; colocalisation sur chaque plateforme.
\sfield{Backtest honnête} Les données de chaque plateforme avec son propre délai~; positions dans la file.
\sfield{Sources} Chaboud, Chiquoine, Hjalmarsson et Vega (2014).
\end{strategyfile}

\begin{strategyfile}{Diffusion déclarée de cours vers des agrégateurs avec last look}\label{strat:hf:fx-market-making:stream}
\sfield{Qui vous paie, et pourquoi} Les clients qui accordent de la valeur à un cours diffusé dans leur propre agrégateur, par segment.
\sfield{Instruments et plateformes} Diffusions vers des banques, des courtiers, des plateformes et des entreprises.
\sfield{Signal} Le \omterm{def:hf:fair-value:fair}{prix juste}~; la toxicité mesurée de chaque segment.
\sfield{Dimensionnement et exécution} Tarifer chaque segment selon ses mark-outs~; n'appliquer le last look que là où il rapporte (le segment toxique dans le modèle).
\sfield{Coûts} Les rejets dont les clients se souviennent~; la couverture.
\sfield{Comment elle disparaît} Les clients qui réorientent leur flux vers la liquidité ferme~; les principes du Code sur le last look et sa divulgation.
\sfield{Horizon, capacité, infrastructure} Quelques millisecondes~; des centaines de diffusions.
\sfield{Backtest honnête} Les rejets tels que les clients les ont vus, et le flux perdu à cause d'eux.
\sfield{Sources} Livre~2, chapitre~15~; le présent chapitre.
\end{strategyfile}

\begin{strategyfile}{Cotation de croisés par synthèse}\label{strat:hf:fx-market-making:cross}
\sfield{Qui vous paie, et pourquoi} Les clients qui négocient des croisés dont les marchés directs sont peu profonds.
\sfield{Instruments et plateformes} Un croisé et ses deux branches passant par une devise commune.
\sfield{Signal} Le \omterm{def:hf:fx-market-making:cross}{croisé synthétique} à partir des cours acheteurs et vendeurs des branches.
\sfield{Dimensionnement et exécution} Coter le croisé à l'intérieur de sa largeur synthétique lorsque le flux peut être internalisé~; couvrir dans les branches ou sur le marché direct,
selon le moins cher.
\sfield{Coûts} Les spreads des deux branches lors de la couverture (2.6 pips sur l'exemple).
\sfield{Comment elle disparaît} Un flux de croisé unilatéral qui ne peut pas être internalisé.
\sfield{Horizon, capacité, infrastructure} Quelques secondes~; des données sur les deux branches et le croisé direct.
\sfield{Backtest honnête} Des prix de branches exécutables au moment de la transaction sur le croisé.
\sfield{Sources} Le présent chapitre.
\end{strategyfile}

\begin{strategyfile}{Teneur de marché électronique sur contrats à terme non livrables}\label{strat:hf:fx-market-making:ndf}
\sfield{Qui vous paie, et pourquoi} Les investisseurs et les entreprises qui couvrent des devises qu'ils ne peuvent pas livrer.
\sfield{Instruments et plateformes} Des contrats à terme non livrables sur des plateformes électroniques.
\sfield{Signal} Les taux onshore lorsqu'ils sont observables, l'historique du fixing, les devises liquides apparentées.
\sfield{Dimensionnement et exécution} Coter des diffusions segmentées~; internaliser entre les échéances et les clients~; couvrir le résiduel dans d'autres contrats à terme.
\sfield{Coûts} Marchés de couverture larges~; risque de fixing.
\sfield{Comment elle disparaît} Les changements de contrôle des capitaux et les litiges sur le fixing.
\sfield{Horizon, capacité, infrastructure} De quelques jours à quelques mois jusqu'au fixing.
\sfield{Backtest honnête} Le fixing tel que publié, avec l'historique de sa source.
\sfield{Sources} Livre~2, chapitre~18.
\end{strategyfile}

\section{Tutoriel~: une centaine de prix à la fois}\label{tut:hf:fx-market-making:tut}

\begin{tutorial}
\textbf{Objectif.} Construire un \omterm{def:hf:fair-value:consolidated}{prix juste consolidé}, le diffuser vers trois segments avec et sans last look, et décider de la part du flux à internaliser.
\textbf{État final~:} le tableau et la \cref{fig:hf:fx-market-making:hedge}.
\begin{tutsteps}
\item \textbf{Le \omterm{def:hf:fair-value:fair}{prix juste}}~: pondérations à variance inverse avec l'obsolescence.
\omcode{code/firm/fxlp/firm_fxlp.py}{34}{38}{Le bruit de chaque plateforme plus la variance du taux depuis sa dernière mise à jour.}\label{lst:hf:fx-market-making:fair}
\item \textbf{Les diffusions} sur le \texttt{firm.lastlook} du livre~2.
\omcode{code/firm/fxlp/firm_fxlp.py}{57}{71}{Les demandes de chaque segment, le contrôle de last look et son analyse des coûts de transaction, pondérés par la part du segment.}\label{lst:hf:fx-market-making:stream}
\item \textbf{Internaliser ou couvrir} avec un skew qui fait venir du flux compensatoire.
\omcode{code/firm/fxlp/firm_fxlp.py}{81}{98}{L'inventaire sous une part de couverture et un skew~; spread, coût de couverture et charge de risque.}\label{lst:hf:fx-market-making:internalise}
\item \textbf{Balayer} la part de couverture pour chaque skew (\texttt{hf\_fx.hedging}).
\end{tutsteps}
\textbf{Pour aller plus loin.} Faire se regrouper dans le temps les clients informés~; ajouter une deuxième paire et internaliser entre elles~; faire tourner les diffusions sur
les plateformes de \texttt{firm.exchsim} avec leurs propres latences.
\end{tutorial}

\section{Construction~: le fournisseur de liquidité de change}\label{bld:hf:fx-market-making:fxlp}

\begin{build}
\bfield{Rôle} Coter le change à partir de nombreuses plateformes, diffuser des prix segmentés avec ou sans last look, et décider de l'internalisation.
\bfield{Interface} \texttt{fair\_price(mids, noise\_bp, ages\_ms, sigma)}, \texttt{venue\_errors},
\texttt{stream(tiers, hold, threshold, policy, sigma, n, seed)}, \texttt{internalise(h, seed, trades, size, sigma\_day, half\_bp, hedge\_bp, lam, skew)}, \texttt{best\_hedge\_share},
\texttt{synthetic\_cross}. Construit sur \texttt{firm.lastlook} (livre~2).
\bfield{Règles} Points de base du taux~; millisecondes~; montants en points de base fois millions.
\bfield{Tests d'acceptation} \texttt{code/firm/fxlp/tests/}~: pondérations à la main et une venue fraîche et précise qui domine~; le prix consolidé bat celui de la
plateforme principale~; aucune transaction rejetée sans last look, toutes les transactions informées rejetées avec le contrôle asymétrique~; couvrir tout ne laisse aucun risque d'inventaire~; un skew plus fort
réduit l'inventaire~; le \omterm{def:hf:fx-market-making:cross}{croisé synthétique} à la main, direct et inversé.
\bfield{Pour aller plus loin} Flux informé regroupé~; plusieurs paires internalisées ensemble~; diffusions sur \texttt{firm.exchsim}.
\end{build}

\begin{omsources}
\item A.~P.~Chaboud, B.~Chiquoine, E.~Hjalmarsson, C.~Vega, Rise of the machines: algorithmic trading in the foreign exchange market, \emph{Journal
of Finance} 69(5), 2014, 2045--2084.
\item Bank for International Settlements, Triennial Central Bank Survey, preliminary results, press release, 30 September 2025.
\end{omsources}

\section{Exercices}

\begin{exercise}[$\star$]\label{exo:hf:fx-market-making:1}
Deux plateformes affichent 1.08500 et 1.08503 avec un bruit de 0.3 et 0.6 point de base, toutes deux fraîches. Quel est le \omterm{def:hf:fair-value:fair}{prix juste}~?
\end{exercise}

\begin{exercise}[$\star$]\label{exo:hf:fx-market-making:2}
Construire l'euro-yen à partir de l'euro-dollar 1.0850/1.0851 et du dollar-yen 150.20/150.21. Quelle est sa largeur en pips~?
\end{exercise}

\begin{exercise}[$\star$]\label{exo:hf:fx-market-making:3}
À 0.4 point de base sur 20~milliards de dollars par jour, que rapporte le flux avant couverture~? Que coûte une couverture intégrale à 0.25~?
\end{exercise}

\begin{exercise}[$\star\star$]\label{exo:hf:fx-market-making:4}
Pourquoi le last look fait-il augmenter la capture du palier toxique et diminuer celle des autres~?
\end{exercise}

\begin{exercise}[$\star\star$]\label{exo:hf:fx-market-making:5}
Pourquoi la meilleure part de couverture est-elle presque toujours proche de 0 ou de 95~\%~?
\end{exercise}

\begin{exercise}[$\star\star$]\label{exo:hf:fx-market-making:6}
Que Chaboud et ses coauteurs ont-ils constaté au sujet de celui qui a retiré l'arbitrage triangulaire~?
\end{exercise}

\begin{exercise}[$\star\star\star$]\label{exo:hf:fx-market-making:7}
\emph{Programmation.} Relancer \texttt{firm.fxlp.stream} avec une attente de 20 millisecondes au lieu de 5. Que deviennent les rejets des clients non informés du palier~3~?
\end{exercise}

\begin{exercise}[$\star\star\star$]\label{exo:hf:fx-market-making:8}
\emph{Trouver la faille.} «~Notre taux d'internalisation est de 90~\%, donc nos coûts de couverture représentent un dixième de ceux de nos concurrents~; nous sommes le fournisseur le moins cher.~»
\end{exercise}

\section{Problème~: Cent prix d'un coup}

\begin{problem}[{Problème du week-end --- cent prix d'un coup}]\label{pb:hf:fx-market-making:1}
Un fournisseur de liquidité non bancaire cote l'euro-dollar sur trois plateformes, le diffuse vers trois paliers et décide de ce qu'il couvre.

\medskip\noindent\textbf{Partie I --- Le marché.}
\begin{enumerate}
\item Résumer l'encadré daté.
\item Que Chaboud et ses coauteurs ont-ils constaté au sujet du trading algorithmique sur le marché des changes~?
\item Comment le \omterm{def:hf:fair-value:fair}{prix juste} est-il construit, et de combien est-il meilleur que celui de la plateforme principale~?
\item Définir un \omterm{def:hf:fx-market-making:cross}{croisé synthétique} et donner la largeur de l'exemple.
\end{enumerate}

\medskip\noindent\textbf{Partie II --- Les flux.}
\begin{enumerate}[resume]
\item Décrire les paliers et le contrôle de last look.
\item Donner la capture et les rejets sans last look et avec le contrôle asymétrique, par palier.
\item Pourquoi le contrôle symétrique est-il moins favorable au fournisseur~?
\item Que coûte le last look aux clients~?
\end{enumerate}

\medskip\noindent\textbf{Partie III --- L'inventaire.}
\begin{enumerate}[resume]
\item Décrire le modèle d'internalisation et le skew.
\item Donner la meilleure part de couverture pour des demi-vies de 1\,150, 870, 770 et 690 transactions.
\item Que rapporte l'internalisation pour une demi-vie de 350 transactions~?
\item Pourquoi le basculement est-il brutal~?
\end{enumerate}

\medskip\noindent\textbf{Partie IV --- Le verdict.}
\begin{enumerate}[resume]
\item Énoncer le \emph{résultat nommé}~: la capture et le taux de rejet du fournisseur de liquidité avec et sans last look, et la part de couverture
à partir de laquelle l'internalisation cesse d'être rentable.
\item Dire auxquels des paliers accorder le last look, et quel prix ferme afficher.
\item Dire comment mesurer la demi-vie de l'inventaire sur les flux réels.
\item Dire ce qui distingue les croisés et les marchés émergents.
\item Quelle fiche de stratégie est la plus exposée à la réglementation~?
\item Que gagne un fournisseur disposant de flux sur de nombreuses paires~?
\item Comment les flux informés regroupés modifieraient-ils les réponses~?
\item En une phrase~: que vend un fournisseur de liquidité de change~?
\end{enumerate}
\end{problem}

\section{Questions d'entretien}

\begin{interviewq}[$\star$ \iqroles{trader}]\label{iq:hf:fx-market-making:1}
Un client prend votre cours vendeur euro-dollar pour 50~millions de dollars. Couvrez-vous maintenant ou attendez-vous~? Qu'est-ce que vous regardez~?
\end{interviewq}

\begin{interviewq}[$\star\star$ \iqroles{chercheur}]\label{iq:hf:fx-market-making:2}
Comment mesureriez-vous la toxicité de chaque client, et que feriez-vous de cette mesure~?
\end{interviewq}

\begin{interviewq}[$\star\star$ \iqroles{développeur}]\label{iq:hf:fx-market-making:3}
Concevez un système qui diffuse 200 prix par palier sur 30 paires et les met à jour dans les 50 microsecondes qui suivent un tick de plateforme.
\end{interviewq}

\begin{interviewq}[$\star\star$ \iqroles{risque}]\label{iq:hf:fx-market-making:4}
Votre taux d'internalisation est passé de 70~\% à 95~\% en un mois. Qu'est-ce qui pourrait l'expliquer, et qu'est-ce qui pourrait mal tourner~?
\end{interviewq}

\begin{interviewq}[$\star\star$ \iqroles{trader}]\label{iq:hf:fx-market-making:5}
Expliquez pourquoi un contrôle de last look symétrique est plus équitable pour les clients qu'un contrôle asymétrique, et ce qu'il coûte au fournisseur.
\end{interviewq}

\begin{interviewq}[$\star\star\star$ \iqroles{chercheur}]\label{iq:hf:fx-market-making:6}
Avec une fourchette $S$, un coût de couverture $hC$ et une charge de risque $\lambda\,v(h)$, trouvez la part de couverture optimale lorsque $v(h)=(1-h)^2V$ et lorsque
$v(h)=(1-h)V'$.
\end{interviewq}
